\scanpage{156}
% archive foliation: 155.
Considérons d'abord \unsure{par} $B\otimes\mathbb F_p$ sur
$A\otimes\mathbb F_p=\mathbb F_p[T]$, donné
par l'équation
\[
-\,S^2-P(T)=0\quad \deleted{\text{en car.}}\quad\text{sur }\mathbb F_p.
\]
\inserted{(supposons $p\ne2$)}
\unsure{Cette} extension est \uline{gén}$^t$ \deleted{\unsure{non ramifié}}
\inserted{étale} ssi $P(T)\not\equiv0\ p$.
Cela montre déjà que \deleted{\unsure{donc}} l'on doit avoir
dans (1) $q\in\{1,2\}$, et dans le cas où $B_{(p)}$ est aussi
normal \inserted{(p.ex. si $\unsure{q|p}=2$)} et si on veut donc que $B\otimes\mathbb F_2$ soit gén.
étale sur $A\otimes\mathbb F_2$, on voit que l'on doit avoir $q=1$.
Mais l'équation
\[
S^2-P(T)(=S^2-\varepsilon Q(T))=0
\]
en car $2$ n'est jamais gén. étale! Donc \deleted{il faut} pour
avoir qque chose gén. étale en \uline{$\forall$ car}, il faut que
$B_{(p)}$ soit non normal \unsure{i.e.} \inserted{que}
\[
P=\underset{\underset{\{\pm1\}}{\uparrow}}{\varepsilon}
\deleted{\underset{\unsure{\{1,2\}}}{q}}Q
\]
soit un carré
mod $4$. \uline{NB} Pour $Q$ fixé, ceci fixe le signe $\varepsilon$ sans
ambiguité, car $-1$ n'est pas un carré mod $4$!). Cette
condition est aussi suffisante\ldots.

\uline{Exemple} (\unsure{lié} aux modules des courbes elliptiques). Soit
$Q(T)=T(T-\lambda)$, avec $\lambda\in\mathbb Z\setminus\{0\}$
(dans le cas qui m'intéresse, \deleted{$\unsure{\lambda}$}
$\lambda=2^\alpha3^\beta$ l'invariant modulaire exceptionnel). On a
\[
Q(T)\equiv T^2\unsure{\equiv}\ (2)\quad\Longleftrightarrow\quad
\lambda\equiv0\ (4).
\]
On trouve
un revêtement étale de
\[
\mathbb E^1_{\mathbb Z}\setminus\{0,\lambda\}
=\operatorname{Spec}\mathbb Z[T]_{T(T-\lambda)}
\]
et $=$ de
\[
\mathbb P^1_{\mathbb Z}\setminus\{0,\lambda\}
=\operatorname{Spec}\mathbb Z[T']_{\lambda T'-1}
\]
en normalisant $\mathbb Z[T']_{(\lambda T'-1)}$ dans
\[
\mathbb Z[T']_{(\lambda T'-1)}/(\delta^2-(\lambda T'-1))
\]
\margin{\unsure{où l'on} $T'=1/T$}
